\chapter{Experimental Methodology}
\label{chap:methodology}

\section{Research Questions}

The evaluation is organised around five research questions.

RQ1: How much does storage backend choice affect checkpoint loading time for real LLM checkpoints?

RQ2: How does checkpoint layout affect achievable performance and backend sensitivity?

RQ3: How much overhead is introduced by the Python integration layer?

RQ4: To what extent do backend gains survive integration into a wider inference stack such as vLLM?

RQ5: What engineering tradeoffs and limitations arise when building a practical checkpoint loading system?

\section{Hardware and Software Environment}

All experiments are conducted on a dedicated Linux system. The hardware specifications are summarised in Table~\ref{tab:environment}.

\begin{table}[h]
\caption{Experimental environment.}
\label{tab:environment}
\centering
\begin{tabular}{ll}
\toprule
Component & Description \\
\midrule
CPU & Intel Core i9-12900K, 16 cores (8P+8E), 24 threads \\
RAM & 32 GB DDR4-3200 \\
Storage & Samsung 970 EVO Plus 500 GB NVMe SSD \\
Filesystem & ext4, default mount options, noatime \\
Kernel & Linux 6.8.1 \\
Python & 3.11 \\
Rust & 1.77 (io\_uring support enabled) \\
\bottomrule
\end{tabular}
\end{table}

The NVMe device is dedicated to benchmark runs and is not shared with the system root filesystem or other workloads. This avoids interference from background processes during measurements. Swap is enabled but is not expected to contribute during benchmark runs given that the checkpoints fit comfortably in RAM.

The system runs Ubuntu 24.04 with the HWE kernel. io\_uring is available and enabled by default on this kernel version.

\section{Backend Configuration}

Each backend is configured to reflect its intended use. The synchronous backend issues blocking \texttt{read()} calls from a thread pool sized to the number of backend threads. The memory-mapped backend opens files with \texttt{mmap} and accesses tensors through page faults without any application-level threading. The io\_uring backend uses fixed pre-registered buffers and batched submission from a single polling thread. Backend configuration is summarised in Table~\ref{tab:backend_config}.

\begin{table}[h]
\caption{Backend configuration by thread model, queue depth, and buffer strategy.}
\label{tab:backend_config}
\centering
\begin{tabular}{llll}
\toprule
Backend & Thread model & Queue depth & Buffer strategy \\
\midrule
Sync & thread pool (4 threads) & 4 & heap-allocated per read \\
Mmap & single-threaded & N/A (fault-driven) & mmap, kernel-managed \\
io\_uring & single-threaded batched & up to 256 & fixed pre-registered \\
\bottomrule
\end{tabular}
\end{table}

The synchronous backend uses 4 threads in the thread pool. This is a conservative setting that reflects typical deployment configurations where the number of loader threads is set equal to the number of model shards. The io\_uring backend is configured with a maximum batch size of 32 requests, which is sufficient to keep the device queue populated on this hardware while limiting the memory footprint of the submission ring.

\section{Models and Checkpoints}

Two real checkpoints are used for headline experiments. The smaller model is GPT-2 (124 million parameters, approximately 550 MB in fp16), downloaded from HuggingFace Hub under the \texttt{gpt2} identifier. The larger model is LLaMA 3.2 1B (1 billion parameters, approximately 2 GB in fp16), downloaded under the \texttt{Llama-3.2-1B} identifier.

GPT-2 is used for development and sanity checks because its small size allows rapid iteration and because its architecture is widely understood. LLaMA 3.2 1B is used for headline results because it is a genuinely modern model with realistic memory access patterns.

Checkpoints are downloaded on first use by the benchmark fixture and cached locally in \texttt{/tmp/tensorstore\_checkpoints}. The cache is cleared between cold-cache measurement runs using the cache eviction procedure described below.

Synthetic GPT-2-like checkpoints are used for controlled reproducibility testing. These checkpoints are generated by the fixture system in \texttt{benchmarks/fixtures.py} and have identical tensor shapes and dtypes to the real GPT-2 model but with randomly initialised weight values. Randomly initialised weights are used because actual model weights are not needed for measuring I/O performance; what matters is the byte layout and access pattern.

SafeTensors format is used for single-file experiments. For partitioned experiments, the ServerlessLLM-style layout is emulated by splitting a single SafeTensors file into multiple partition files of equal size, following the partitioning strategy described in the ServerlessLLM paper~\cite{FU:2024}.

\section{Cache Eviction Procedure}

Cold-cache runs are the primary measurement condition. Before each cold run, the following procedure is applied:

\begin{enumerate}
\item The benchmark process exits and the Python interpreter is restarted. This ensures that no Python-level caches retain checkpoint data.
\item The page cache is cleared by calling \texttt{posix\_fadvise(fd, 0, 0, POSIX\_FADV\_DONTNEED)} on each checkpoint file. This advises the kernel that the specified pages are no longer needed and should be evicted from the page cache.
\item The filesystem metadata cache is cleared using \texttt{echo 3 > /proc/sys/vm/drop\_caches}. This requires root privileges and is only used for cold-cache runs.
\item The benchmark process is started and the measurement run begins.
\end{enumerate}

Cache eviction is verified by running \texttt{iostat -x 1} in parallel with the benchmark. A successful cold-cache run shows non-zero \texttt{\%util} (device utilisation) on the NVMe device and shows actual read operations in the \texttt{await} (average wait time) column. A failed eviction (where the page cache was not cleared) shows near-zero device utilisation and very low await values.

Warm-cache runs follow the same procedure but skip step 2 and step 3, allowing the page cache to retain checkpoint data between runs.

\section{Measurement Procedure}

Each measurement run follows this sequence:

\begin{enumerate}
\item Cache eviction (cold runs only).
\item Warm-up run: one iteration of the benchmark, results discarded.
\item Measurement runs: five iterations, results averaged.
\item Cache re-eviction (cold runs only, before next run).
\end{enumerate}

Five measurement iterations are used because pilot experiments showed that variance across runs was low (typically below 2 percent coefficient of variation) after cache eviction, and five iterations provide a stable mean without excessive runtime.

Each iteration reports the wall-clock time from the first read request to the last tensor being available in memory. The median of the five iterations is reported as the primary result, with the min and max reported as a range.

Warm-up runs are performed to ensure that any one-time initialisation costs (such as JIT compilation of Python paths, dynamic linking of Rust libraries, or one-time page faults for code pages) do not contaminate the measurement.

\section{Metrics}

The primary metric is end-to-end checkpoint load time, measured as wall-clock time from the first read request to the last tensor being available in memory. This metric directly reflects user-observable cold-start latency.

Secondary metrics are:

\begin{itemize}
\item Effective throughput: total bytes read divided by load time, expressed in GB/s.
\item Read IOPS: number of completed read operations divided by load time.
\item Per-read latency percentiles (p50, p95, p99): measured at the application level as the interval between submitting a request and receiving its data.
\end{itemize}

Effective throughput and IOPS are derived from the same wall-clock measurement and provide a device-agnostic way to compare backends across different hardware. Per-read latency percentiles are collected by instrumenting the backend to record the submission and completion time of each individual read operation.

\section{Benchmark Kinds}

Three benchmark kinds are used.

Load-only measures the time to load model weights into memory without performing any inference. This isolates the loading operation from the rest of the model initialisation process and allows pure comparison of I/O performance.

TTFT (time-to-first-token) measures the interval from when a request is submitted to when the first generated token is returned. This is measured in the vLLM integration benchmarks and includes both loading time and the vLLM-specific preparation steps.

Steady-state decode measures per-token latency during continued generation after the model is loaded and warmed up. This benchmark is used to verify that the loading backend does not introduce overhead during the decode phase.

\section{Fairness Controls}

Request traces are generated once per checkpoint and reused across backends without modification. This ensures that the same logical workload is evaluated in each case. Backends are not permitted to alter request boundaries, ordering, or grouping.

The evaluation harness is entirely separate from the loader code. No timing instrumentation is added to the loader itself; instead, wall-clock time is measured at the harness level using Python's \texttt{time.perf\_counter()}, which provides nanosecond-resolution monotonic timing.

Backends that support configurable parameters (such as thread pool size or batch size) are configured identically across comparison runs. The configuration used for each backend is reported in Table~\ref{tab:backend_config}.

\section{Statistical Analysis}

Results are reported as the mean of five runs with standard deviation. The coefficient of variation (CV, standard deviation divided by mean) is used to assess result stability. A CV below 5 percent is considered stable; a CV above 10 percent triggers investigation of the source of variance.

Where comparisons between backends are made, the difference is reported as a percentage improvement over the synchronous baseline, calculated as $(t_{baseline} - t_{backend}) / t_{baseline} \times 100$.

Significance is not formally assessed because the experiments are deterministic (the same checkpoint always produces the same access pattern) and variance is low after cache eviction. The reporting of means, standard deviations, and ranges is considered sufficient for characterising the results.

\section{Threats to Validity}

Several threats to validity apply.

Internal validity is threatened by page cache contamination. This is addressed by the cache eviction procedure and iostat verification described above.

External validity is threatened by the single hardware configuration. All experiments are conducted on one machine with one NVMe model; absolute throughput and latency values may differ on other devices, particularly on older storage hardware or network-attached storage. The qualitative ordering of backends (which backend performs best under which conditions) is expected to generalise to other NVMe devices, but this claim is not directly verified.

Construct validity is threatened by the use of synthetic checkpoints for some experiments. Synthetic checkpoints have identical access patterns to real checkpoints but may differ in ways that affect performance (for example, filesystem compression or deduplication). This is mitigated by using real HuggingFace checkpoints for headline results.

The vLLM integration benchmarks measure loading time in isolation and do not include end-to-end serving performance under concurrent requests. The interaction between loading backend choice and multi-request serving performance is left to future work.

These limitations are discussed further in Chapter~\ref{chap:discussion}.
